% ============================================================
% RESEARCH STATEMENT
% Standalone document. Compile twice with pdfLaTeX.
% ============================================================

\documentclass[11pt,letterpaper]{article}

\usepackage[
    margin=0.75in,
    headheight=14pt,
    headsep=14pt,
    footskip=22pt
]{geometry}

\usepackage[T1]{fontenc}
\usepackage[utf8]{inputenc}
\usepackage{newtxtext}
\usepackage{amsmath}
\usepackage{newtxmath}
\usepackage{microtype}
\usepackage{xcolor}
\usepackage{titlesec}
\usepackage{fancyhdr}
\usepackage[hidelinks]{hyperref}

% ============================================================
% COLORS
% ============================================================

\definecolor{Accent}{HTML}{007F6B}
\definecolor{Body}{HTML}{222222}
\definecolor{Muted}{HTML}{555555}

% ============================================================
% SPACING AND TYPOGRAPHY
% ============================================================

\setlength{\parindent}{0pt}
\setlength{\parskip}{3.5pt}
\linespread{1.0}

\raggedbottom
\emergencystretch=1.5em
\clubpenalty=10000
\widowpenalty=10000

\titleformat{\section}
    {\large\bfseries\color{Accent}}
    {}
    {0pt}
    {}

\titlespacing*{\section}
    {0pt}
    {8pt}
    {3pt}

\titleformat{\paragraph}[runin]
    {\normalsize\bfseries}
    {}
    {0pt}
    {}
    [.]

\titlespacing*{\paragraph}
    {0pt}
    {5pt}
    {0.5em}

% ============================================================
% HEADERS AND FOOTERS
% ============================================================

\pagestyle{fancy}
\fancyhf{}

\fancyhead[L]
    {\small\color{Muted}Mohammad Alipour-Vaezi}

\fancyhead[R]
    {\small\color{Muted}Research Statement}

\fancyfoot[C]
    {\small\color{Muted}\thepage}

\renewcommand{\headrulewidth}{0.3pt}
\renewcommand{\footrulewidth}{0pt}

\fancypagestyle{firstpage}{%
    \fancyhf{}
    \fancyfoot[C]{\small\color{Muted}\thepage}
    \renewcommand{\headrulewidth}{0pt}
}

\hypersetup{
    pdftitle={
        Research Statement:
        Reliable Learning and Decision-Making Under Risk
    },
    pdfauthor={Mohammad Alipour-Vaezi}
}

% ============================================================
% DOCUMENT
% ============================================================

\begin{document}

\color{Body}
\thispagestyle{firstpage}

{\LARGE\bfseries Research Statement}
\par
\vspace{3pt}

{\large Mohammad Alipour-Vaezi}
\par

{\small\color{Muted}
Industrial and Systems Engineering \textbar{} Virginia Tech}
\par
\vspace{5pt}

{\color{Accent}\hrule height 0.6pt}
\vspace{6pt}

{\large\bfseries
Reliable Learning and Decision-Making Under Risk}
\par

My research develops methods for \textbf{learning to make consequential
decisions from imperfect information}. The central question is not only
which decision performs well on average, but which decision remains
defensible when outcomes are unfavorable, evidence is limited, and
resources are constrained. A healthcare system must consider both patient
outcomes and service capacity; an artificial intelligence agent must
decide not only what action to take, but whether its information is
reliable enough to act. I connect these problems through reinforcement
learning, data-driven optimization, and multi-attribute decision-making.

My long-term goal is to build a research program in which \textbf{risk
preferences, the credibility of evidence, and operational constraints
are modeled together}. The program has three connected directions:
foundations for risk-sensitive sequential learning; causal policy learning
for personalized, capacity-constrained healthcare; and risk-aware
information management for language-model agents. The first develops the
mathematical foundation, the second addresses observational evidence, and
the third makes information acquisition and retention part of the decision
itself.

\section{Research foundation and current work}

\paragraph{Quantile objectives for sequential learning}
My AISTATS 2026 paper, \emph{Optimistic Reinforcement Learning with Quantile
Objectives}, and my job-market paper, \emph{Risk-Sensitive Reinforcement
Learning with Smoothed Quantile Objectives}, form a connected line of work
on learning policies for distributional objectives~\cite{aistats,smoothed}.
The job-market paper, currently under review at \emph{Operations Research},
addresses the instability of point quantiles when transition probabilities
must be estimated. With my coauthors, I develop UCB--BQRL, an optimistic
model-based learning method using a lower-buffered quantile criterion, and
EVI--BQ, an exact planning procedure. The manuscript establishes a
high-probability regret bound, an information-theoretic lower bound, and
computational hardness results for exact quantile evaluation. These
results motivate my next question: when does statistical learnability
translate into a computationally practical method?

My ongoing manuscript, \emph{Deep Reinforcement Learning with Buffered
Quantile Objectives}, extends this line toward neural function
approximation~\cite{deep}. It develops model-free learning of conditional
return quantiles, buffered action scoring, and uncertainty-guided
exploration. This provides a starting point for characterizing when an
approximate learned distribution supports a reliable policy.

\paragraph{Healthcare and multi-attribute decisions}
My applied research includes published work on emergency-department
prioritization and queueing~\cite{ed}, as well as studies of health outcomes.
Current projects examine multisystem outcomes associated with glucagon-like
peptide-1 receptor agonist use, personalized dosage adjustment for
incretin-based therapies, and an Augmented Numeric Rating Scale for pain
measurement. They motivate longitudinal decisions that account for
uncertainty, patient preferences, and limited resources. My work on
construction portfolio optimization and motion-picture portfolio
management further connects multi-attribute evaluation with resource
allocation; the latter includes the use of a language model as an expert
input~\cite{portfolio}.

\paragraph{Information as a decision resource}
My current applied-science internship at Amazon focuses on risk-sensitive
agentic AI. My ongoing work on task-adaptive quantile utility for memory
management distinguishes whether stored information is relevant, whether
its content is trustworthy, and whether using it improves the task
outcome. This motivates a broader academic question: when should an
adaptive system reuse information, acquire new evidence, or defer a
decision? My academic program will use publishable research and
appropriately authorized data rather than depend on proprietary
infrastructure.

\section{A common framework: risk, evidence, and constraints}

I will study policies that choose operational actions together with
information actions, such as requesting an observation, verifying a
memory, or seeking human input. A representative target is
\begin{equation}
\begin{aligned}
\max_{\pi\in\Pi}\quad
    &\inf_{P\in\mathcal{P}_n}
      \mathcal{B}^{P}_{\tau,\beta}(G^{\pi})
    \\
\text{subject to}\quad
    &\sup_{P\in\mathcal{P}_n}P(F^{\pi})\leq\alpha,
    \qquad
    \sup_{P\in\mathcal{P}_n}
    \mathbb{E}_{P}[C^{\pi}]\leq B.
\end{aligned}
\label{eq:agenda}
\end{equation}
Here, $\Pi$ is a specified policy class; $G^{\pi}$ is cumulative utility,
with larger values preferred; $F^{\pi}$ is an unacceptable-outcome event;
and $C^{\pi}$ is resource consumption. The tolerable failure probability
and resource budget are $\alpha$ and $B$. The set $\mathcal{P}_n$
represents models supported by $n$ observations and explicitly stated
sensitivity assumptions. Under model $P$, the lower-buffered criterion is
\[
\mathcal{B}^{P}_{\tau,\beta}(G^{\pi})
    =
    \frac{1}{\beta}
    \int_{\tau-\beta}^{\tau}
    Q^{P}_{u}(G^{\pi})\,\mathrm{d}u,
\qquad
0<\beta\leq\tau<1,
\]
where $Q^{P}_{u}$ denotes the $u$-quantile. The decision maker selects the
target quantile $\tau$; $\beta$ controls the smoothing interval. Neither
is the statistical confidence level of an estimated guarantee. A separate
failure constraint prevents favorable performance near one quantile from
concealing unacceptable events elsewhere in the distribution.

Equation~\eqref{eq:agenda} is a research target. Meaningful guarantees
require justified model coverage and feasible constraints. I will
characterize these conditions and develop rules that retain a baseline
or request evidence when improvement cannot be supported. The three
directions address different limitations: computation, causal
identification, and information access.

\section{Direction 1: Computationally practical risk-sensitive learning}

\textbf{How much data and computation are needed to improve a policy at a
specified level of risk?} Building on my quantile-learning work, I will
develop algorithms whose guarantees account jointly for statistical
uncertainty, planning error, and the approximation introduced by
smoothing. My initial setting will be finite-horizon problems with
structured dynamics and clearly specified policy classes. I will identify
which properties of the return distribution determine the difficulty of
learning.

The first project will develop approximate planning procedures with
explicit error certificates and study how their errors accumulate during
learning. I will investigate adaptive allocation of computation near
decision-relevant quantiles rather than uniformly across the entire
return distribution. The targets are upper and lower bounds exposing the
tradeoff among sample size, computation, and policy quality, together with
structural assumptions that permit efficient approximation.

A second project will develop \emph{risk-conditioned policies} that support
a range of stakeholder preferences without retraining from scratch for
each quantile. I will investigate uniform performance bounds across a
specified range of quantile levels, the cost of sharing information
across those levels, and the effect of preference uncertainty. When
preferences change during a trajectory, I will distinguish an initially
committed objective from a dynamically updated one and study the augmented
state needed for each. The aim is to establish when adaptation preserves
the intended distributional objective.

A third project will examine learning under changing environments and
limited feedback. The emphasis will be on retaining performance relative
to a baseline while identifying when additional exploration is justified.
Evaluation will compare mean return, target-quantile return, failure
frequency, sample use, and computation against risk-neutral,
risk-sensitive, and robust baselines. Sequential allocation and
service-system experiments will reveal when reliability requires more
data, more computation, or a sacrifice in average performance.

\section{Direction 2: Causal and resource-aware personalized decisions}

\textbf{When can historical data justify a change in a sequential treatment
or service policy?} My healthcare projects motivate a program that connects
distributional causal inference with risk-sensitive optimization. I will
begin with longitudinal treatment and follow-up decisions, using
incretin-based therapies as a motivating application rather than defining
the entire program around one drug class. Policies must address both
heterogeneous outcomes and implementation resources.

The first project will estimate outcome distributions under candidate
intervention policies. I will make consistency, sequential exchangeability,
and treatment overlap explicit, and study estimators combining outcome and
treatment-assignment models. The target is a distribution under a policy,
not an unsupported claim about the distribution of individual treatment
effects. When confounding cannot be ruled out or data provide weak support
for particular actions, I will develop sensitivity-indexed bounds and
restricted policy classes. This will identify where evidence does and
does not support personalization.

The second project will translate those estimates into conservative policy
improvement. I will investigate joint uncertainty bounds for a candidate
policy and its baseline, with held-out evaluation or appropriately uniform
guarantees to account for policy selection. Patient-valued outcomes,
adverse-event thresholds, and treatment burden will remain distinguishable
rather than being concealed in an arbitrary aggregate score. My
multi-attribute decision-making work motivates preference elicitation and
sensitivity analysis over plausible priorities. I will investigate when
further data have greater decision value than a more complex algorithm.

The third project will move from patient-level recommendations to
\emph{coupled decisions across a care system}. Appointment capacity,
monitoring resources, and treatment budgets connect decisions that might
otherwise be optimized separately. I will study allocation rules that
balance population benefit with explicitly defined protection for groups
of patients, including the additional uncertainty created by small
subgroup samples. The operations question is how to allocate capacity
when response distributions and the consequences of delay are uncertain.

I will begin with simulation and semi-synthetic experiments where policy
consequences are known, followed by retrospective evaluation using
appropriately accessible longitudinal data. Clinical collaborations will
be developed for problem definition, data access, and external validation.
Prospective use would require separate approval and validation;
retrospective policy estimates would not be presented as demonstrated
clinical benefit.

\section{Direction 3: Risk-aware information management for AI agents}

\textbf{What should an agent remember, verify, or forget when information
has uncertain downstream value?} I will extend task-adaptive quantile
utility from a memory-selection criterion to a sequential theory of
information management. The central hypothesis is that evaluating
information through its contribution to the distribution of task outcomes
will produce better reliability--cost tradeoffs than evaluating relevance
or average usefulness alone.

The first project will estimate the effect of information-use decisions.
In controlled environments, randomized inclusion, exclusion, and
verification of memories will support comparisons between alternative
memory policies. I will study delayed feedback and interactions among
memories, including cases where two individually useful memories conflict
when used together. The inferential target will be the change in the
outcome distribution induced by an information policy, with separate
checks on factual trustworthiness.

The second project will integrate retrieval, verification, updating,
retirement, and human escalation into a resource-constrained sequential
model. I will investigate sufficient memory representations and conditions
under which simple verification or retirement rules approximate a more
complex policy. The theoretical challenge is to characterize the value
of information under a distributional objective when storing or discarding
information changes future decisions. This extends Direction 1 to
choosing which uncertainty to reduce.

The third project will study temporal deterioration and transfer across
task classes. Evaluation will separately vary stale information,
conflicting memories, task consequences, and feedback availability.
Comparisons will include relevance-only retrieval, mean-utility selection,
fixed-risk criteria, and task-adaptive policies under matched memory and
computation budgets. Outcomes will include task success, lower-tail
utility, harmful information use, escalation burden, and recovery after
updates. Guarantees will be stated for specified task groups and shift
assumptions; success on one benchmark will not be treated as evidence of
universal reliability.

\section{Doctoral training and a sustainable research program}

The agenda supports \textbf{two initial, scientifically distinct doctoral
trajectories}. A dissertation on \emph{Reliable Sequential Learning Under
Computational and Data Constraints} would connect approximate quantile
planning, learning across risk preferences, and baseline-aware adaptation.
This work can begin with controlled models and reproducible simulations
without requiring clinical or industrial data.

A second dissertation on \emph{Distributional Causal Policy Learning for
Personalized and Capacity-Constrained Care} would connect identification
and estimation of policy-outcome distributions, sensitivity-aware policy
improvement, and allocation under shared resources. Initial theory and
semi-synthetic evaluation can proceed independently of the first
dissertation, incorporating improved planning methods as they become
available. Shared infrastructure would support collaboration while
preserving intellectual ownership.

As the group grows, the information-management direction provides a third
doctoral pathway on \emph{Risk-Aware Memory and Verification for Adaptive
Agents}. Shared software for policy evaluation and stress testing would
also provide entry points for undergraduate and master's research.

\section{External funding and the first five years}

My funding strategy will pair a coherent methodological foundation with
distinct application-driven proposals. For risk-sensitive learning and
information control, I would discuss fit with NSF's
\href
{https://www.nsf.gov/funding/opportunities/iids-intelligent-interactive-dynamic-systems}
{\emph{Intelligent and Interactive Dynamic Systems}}
program, whose scope includes inference, control, reliability, and systems
that adapt to people. Operational extensions concerning manufacturing and
service-system design could support separate proposals to NSF's
\href
{https://www.nsf.gov/funding/opportunities/amdo-advanced-manufacturing-design-operations}
{\emph{Advanced Manufacturing, Design, and Operations}}
program when those applications are central. Each proposal would have
distinct aims and deliverables.

For healthcare, I would develop collaborative NIH proposals aligned with
\href
{https://www.niddk.nih.gov/research-funding/research-programs}
{NIDDK's research interests in diabetes and obesity},
with clinical partners shaping the substantive questions and validation
plan. Retrospective feasibility would precede external validation, with
intervention studies proposed separately. Industry-supported projects
would address publishable questions in agent reliability and evaluation.
Once eligible, I would develop an NSF
\href
{https://www.nsf.gov/funding/opportunities/career-faculty-early-career-development-program}
{\emph{CAREER}}
proposal integrating the foundational research with courses and open
teaching modules on reliable sequential decisions.

In years one and two, I will prioritize the first two doctoral
trajectories, reusable evaluation tools, pilot studies, and initial
methods proposals. In years three and four, I will extend policy classes
and datasets and pursue larger collaborative proposals supported by
preliminary evidence. By year five, the goal is a group with complementary
theoretical and applied contributions and a platform for subsequent
doctoral work. Expansion will be tied to secured resources and feasible
data access.

The lasting contribution I seek is a set of principles for deciding
\textbf{what to do, what evidence to obtain, and when not to act} under
risk. Such principles would connect mathematical advances in
reinforcement learning with operational choices about personalization,
capacity, and automation.

% ============================================================
% SELECTED REFERENCES
% Publication and manuscript statuses follow the supplied CV.
% ============================================================

\renewcommand{\refname}{Selected Research Referenced}

\begin{thebibliography}{9}

\small
\setlength{\itemsep}{2pt}
\setlength{\parskip}{0pt}

\bibitem{aistats}
M.~Alipour-Vaezi, H.~Zhong, K.~L.~Tsui, and S.~Khodadadian.
Optimistic Reinforcement Learning with Quantile Objectives.
\emph{International Conference on Artificial Intelligence and Statistics},
2026.

\bibitem{smoothed}
M.~Alipour-Vaezi, H.~Zhong, and S.~Khodadadian.
Risk-Sensitive Reinforcement Learning with Smoothed Quantile Objectives.
Manuscript under review at \emph{Operations Research}.

\bibitem{deep}
M.~Alipour-Vaezi and S.~Khodadadian.
Deep Reinforcement Learning with Buffered Quantile Objectives.
Manuscript under review at ICLR.

\bibitem{ed}
M.~Alipour-Vaezi, A.~Aghsami, and F.~Jolai.
Prioritizing and queueing the emergency departments' patients using a
novel data-driven decision-making methodology, a real case study.
\emph{Expert Systems with Applications}, 195:116568, 2022.

\bibitem{portfolio}
M.~Alipour-Vaezi and K.~L.~Tsui.
Data-driven portfolio management for motion pictures industry:
A new data-driven optimization methodology using a large language
model as the expert.
\emph{Computers \& Industrial Engineering}, 197:110574, 2024.

\end{thebibliography}

\end{document}